% ======================================================================
%  UPPGIFT 5: UNDRE GRÄNS FÖR SÖKNING
% ======================================================================
\section{5 Undre gränser}

\uppgiftsbild{5}{Undre gräns för sökning i sorterad array}
{bevisa en undre gräns för ett \emph{problem} — och hålla isär den från gränser för algoritmer
och från bästa och värsta fall.}
{\begin{itemize}\setlength{\itemsep}{1pt}
  \item Gränsens riktning ($O$, $\Omega$, $\Theta$) och fallet (bästa, medel, värsta) är två
        olika saker.
  \item Algoritm eller problem: $\exists A\,\forall x$ mot $\forall A\,\exists x$.
  \item Beslutsträdet: löven är möjliga svar, höjden är värsta fallet.
\end{itemize}}

\begin{frame}{Övre och undre gräns är inte bästa och värsta fall}
\centering
\begin{tikzpicture}[x=1cm,y=1cm,
  hd/.style={font=\small\bfseries,text=kthnavy,align=center},
  rl/.style={font=\small,align=left,anchor=west},
  cell/.style={draw=black!25,minimum width=30.5mm,minimum height=9.2mm,inner sep=0pt},
  big/.style={font=\normalsize},
  loose/.style={font=\scriptsize,text=black!50}]
  \useasboundingbox (-0.05,0.55) rectangle (11.95,-3.95);
  \node[hd] at (4.05,0) {bästa fallet\\[-1pt]{\footnotesize\mdseries $T_{\min}(n)$}};
  \node[hd] at (7.2,0) {medelfallet\\[-1pt]{\footnotesize\mdseries $T_{\mathrm{medel}}(n)$}};
  \node[hd] at (10.35,0) {värsta fallet\\[-1pt]{\footnotesize\mdseries $T_{\max}(n)$}};
  \node[rl] at (0,-0.95) {$O$\ \ övre gräns};
  \node[rl] at (0,-1.95) {$\Omega$\ \ undre gräns};
  \node[rl] at (0,-2.95) {$\Theta$\ \ båda};
  \node[cell] at (4.05,-0.95) {};
  \node[cell] at (4.05,-1.95) {};
  \node[cell] at (4.05,-2.95) {};
  \node[cell] at (7.2,-0.95) {};
  \node[cell] at (7.2,-1.95) {};
  \node[cell] at (7.2,-2.95) {};
  \node[cell] at (10.35,-0.95) {};
  \node[cell] at (10.35,-1.95) {};
  \node[cell] at (10.35,-2.95) {};
  \node[cell,fill=coral!25,visible on=<1>] at (10.35,-0.95) {\small $O$ = värsta};
  \node[cell,fill=coral!25,visible on=<1>] at (7.2,-2.95) {\small $\Theta$ = medel};
  \node[cell,fill=coral!25,visible on=<1>] at (4.05,-1.95) {\small $\Omega$ = bästa};
  \node[rotate=12,draw=brick,line width=2pt,text=brick,rounded corners=3pt,font=\Huge\bfseries,inner sep=4pt,fill=white,fill opacity=0.85,text opacity=1,visible on=<1>] at (7.2,-1.95) {Nej.};
  \node[cell,fill=lampyellow!45,draw=kthorange,line width=1pt,visible on=<2->] at (4.05,-0.95) {};
  \node[big,visible on=<2>] at (4.05,-0.95) {$O(n)$};
  \node[big,visible on=<3>] at (4.05,-0.8099999999999999) {$O(n)$};
  \node[loose,visible on=<3>] at (4.05,-1.19) {också $O(n^2)$, $O(n^3)$};
  \node[big,visible on=<2>] at (7.2,-0.95) {$O(n^2)$};
  \node[big,visible on=<3>] at (7.2,-0.8099999999999999) {$O(n^2)$};
  \node[loose,visible on=<3>] at (7.2,-1.19) {också $O(n^3)$};
  \node[big,visible on=<2>] at (10.35,-0.95) {$O(n^2)$};
  \node[big,visible on=<3>] at (10.35,-0.8099999999999999) {$O(n^2)$};
  \node[loose,visible on=<3>] at (10.35,-1.19) {också $O(n^3)$, $O(2^n)$};
  \node[big,visible on=<2>] at (4.05,-1.95) {$\Omega(n)$};
  \node[big,visible on=<3>] at (4.05,-1.81) {$\Omega(n)$};
  \node[loose,visible on=<3>] at (4.05,-2.19) {också $\Omega(1)$};
  \node[big,visible on=<2>] at (7.2,-1.95) {$\Omega(n^2)$};
  \node[big,visible on=<3>] at (7.2,-1.81) {$\Omega(n^2)$};
  \node[loose,visible on=<3>] at (7.2,-2.19) {också $\Omega(n)$};
  \node[cell,fill=lampyellow!45,draw=kthorange,line width=1pt,visible on=<2->] at (10.35,-1.95) {};
  \node[big,visible on=<2>] at (10.35,-1.95) {$\Omega(n^2)$};
  \node[big,visible on=<3>] at (10.35,-1.81) {$\Omega(n^2)$};
  \node[loose,visible on=<3>] at (10.35,-2.19) {också $\Omega(n\log n)$};
  \node[big,visible on=<2>] at (4.05,-2.95) {$\Theta(n)$};
  \node[big,visible on=<3>] at (4.05,-2.81) {$\Theta(n)$};
  \node[loose,visible on=<3>] at (4.05,-3.1900000000000004) {bara den klassen};
  \node[big,visible on=<2>] at (7.2,-2.95) {$\Theta(n^2)$};
  \node[big,visible on=<3>] at (7.2,-2.81) {$\Theta(n^2)$};
  \node[loose,visible on=<3>] at (7.2,-3.1900000000000004) {bara den klassen};
  \node[big,visible on=<2>] at (10.35,-2.95) {$\Theta(n^2)$};
  \node[big,visible on=<3>] at (10.35,-2.81) {$\Theta(n^2)$};
  \node[loose,visible on=<3>] at (10.35,-3.1900000000000004) {bara den klassen};
  \node[visible on=<2->,font=\footnotesize\bfseries,text=kthorange] at (4.05,-3.68) {\raisebox{0.1ex}{$\blacktriangle$} Bästa fallet har ett ordo};
  \node[visible on=<2->,font=\footnotesize\bfseries,text=kthorange] at (10.35,-3.68) {\raisebox{0.1ex}{$\blacktriangle$} Värsta fallet har ett $\Omega$};
\end{tikzpicture}
\vfill
\begin{tcolorbox}[enhanced,halign=left,colback=kthlight!50,colframe=kthsky,boxrule=0.4pt,arc=2pt,left=4pt,right=4pt,top=1pt,bottom=1pt,fontupper=\small,height=1.25cm,valign=center]
\only<1>{En vanlig bild: tre synonymer. $O$ betyder värsta fallet, $\Omega$ bästa, $\Theta$ medel. Den är fel, och den gör resten av kursen svårare än den behöver vara.}\only<2>{Insättningssortering, alla nio rutor. \textbf{Fallet} väljer \emph{vilken funktion} vi pratar om. $O$, $\Omega$, $\Theta$ säger hur \emph{den} växer: uppåt begränsad, nedåt begränsad, eller båda.}\only<3>{$O$ och $\Omega$ får vara lösa och fortfarande sanna. $\Theta$ får inte det. ”Värsta fallet är $O(2^n)$” är sant — och helt ointressant.}
\end{tcolorbox}
\end{frame}


\begin{frame}{Tre funktioner, inte en}
\begin{columns}[T]
\begin{column}{0.53\textwidth}
\centering
{\footnotesize\bfseries\color{kthnavy}Insättningssortering: jämförelser}\par\smallskip
\begin{tikzpicture}[x=0.56cm,y=0.054cm]
  \draw[->,black!60] (0,0) -- (11,0) node[below left,font=\scriptsize,text=black!70] {$n$};
  \draw[->,black!60] (0,0) -- (0,50) node[above right,font=\scriptsize,text=black!70] {jämförelser};
  \draw[black!40] (2,0) -- (2,-0.8) node[below,font=\tiny,text=black!60] {2};
  \draw[black!40] (4,0) -- (4,-0.8) node[below,font=\tiny,text=black!60] {4};
  \draw[black!40] (6,0) -- (6,-0.8) node[below,font=\tiny,text=black!60] {6};
  \draw[black!40] (8,0) -- (8,-0.8) node[below,font=\tiny,text=black!60] {8};
  \draw[black!40] (10,0) -- (10,-0.8) node[below,font=\tiny,text=black!60] {10};
  \draw[black!40] (0,10) -- (-0.15,10) node[left,font=\tiny,text=black!60] {10};
  \draw[black!40] (0,20) -- (-0.15,20) node[left,font=\tiny,text=black!60] {20};
  \draw[black!40] (0,30) -- (-0.15,30) node[left,font=\tiny,text=black!60] {30};
  \draw[black!40] (0,40) -- (-0.15,40) node[left,font=\tiny,text=black!60] {40};
  \draw[black!15,line width=3.2pt] (1,0) -- (1,0.0);
  \draw[black!15,line width=3.2pt] (2,1) -- (2,1.0);
  \draw[black!15,line width=3.2pt] (3,2) -- (3,3.0);
  \draw[black!15,line width=3.2pt] (4,3) -- (4,6.0);
  \draw[black!15,line width=3.2pt] (5,4) -- (5,10.0);
  \draw[black!15,line width=3.2pt] (6,5) -- (6,15.0);
  \draw[black!15,line width=3.2pt] (7,6) -- (7,21.0);
  \draw[black!15,line width=3.2pt] (8,7) -- (8,28.0);
  \draw[black!15,line width=3.2pt] (9,8) -- (9,36.0);
  \draw[black!15,line width=3.2pt] (10,9) -- (10,45.0);
  \draw[brick,line width=1pt] (1,0.000) -- (2,1.000) -- (3,3.000) -- (4,6.000) -- (5,10.000) -- (6,15.000) -- (7,21.000) -- (8,28.000) -- (9,36.000) -- (10,45.000);
  \fill[brick] (1,0.000) circle (1.6pt);
  \fill[brick] (2,1.000) circle (1.6pt);
  \fill[brick] (3,3.000) circle (1.6pt);
  \fill[brick] (4,6.000) circle (1.6pt);
  \fill[brick] (5,10.000) circle (1.6pt);
  \fill[brick] (6,15.000) circle (1.6pt);
  \fill[brick] (7,21.000) circle (1.6pt);
  \fill[brick] (8,28.000) circle (1.6pt);
  \fill[brick] (9,36.000) circle (1.6pt);
  \fill[brick] (10,45.000) circle (1.6pt);
  \draw[kthorange,line width=1pt] (1,0.000) -- (2,1.000) -- (3,2.667) -- (4,4.917) -- (5,7.717) -- (6,11.050) -- (7,14.907) -- (8,19.282) -- (9,24.171) -- (10,29.571);
  \fill[kthorange] (1,0.000) circle (1.6pt);
  \fill[kthorange] (2,1.000) circle (1.6pt);
  \fill[kthorange] (3,2.667) circle (1.6pt);
  \fill[kthorange] (4,4.917) circle (1.6pt);
  \fill[kthorange] (5,7.717) circle (1.6pt);
  \fill[kthorange] (6,11.050) circle (1.6pt);
  \fill[kthorange] (7,14.907) circle (1.6pt);
  \fill[kthorange] (8,19.282) circle (1.6pt);
  \fill[kthorange] (9,24.171) circle (1.6pt);
  \fill[kthorange] (10,29.571) circle (1.6pt);
  \draw[kthgrass,line width=1pt] (1,0.000) -- (2,1.000) -- (3,2.000) -- (4,3.000) -- (5,4.000) -- (6,5.000) -- (7,6.000) -- (8,7.000) -- (9,8.000) -- (10,9.000);
  \fill[kthgrass] (1,0.000) circle (1.6pt);
  \fill[kthgrass] (2,1.000) circle (1.6pt);
  \fill[kthgrass] (3,2.000) circle (1.6pt);
  \fill[kthgrass] (4,3.000) circle (1.6pt);
  \fill[kthgrass] (5,4.000) circle (1.6pt);
  \fill[kthgrass] (6,5.000) circle (1.6pt);
  \fill[kthgrass] (7,6.000) circle (1.6pt);
  \fill[kthgrass] (8,7.000) circle (1.6pt);
  \fill[kthgrass] (9,8.000) circle (1.6pt);
  \fill[kthgrass] (10,9.000) circle (1.6pt);
  \node[font=\scriptsize\bfseries,text=brick,anchor=west] at (10.2,45) {$T_{\max}$};
  \node[font=\scriptsize\bfseries,text=kthorange,anchor=west] at (10.2,29.57) {$T_{\mathrm{medel}}$};
  \node[font=\scriptsize\bfseries,text=kthgrass,anchor=west] at (10.2,9) {$T_{\min}$};
\end{tikzpicture}

\end{column}
\begin{column}{0.44\textwidth}
\footnotesize
För varje $n$ finns $n!$ indata (grå staplar). Varje \emph{fall} plockar ut en funktion av $n$:
\begin{itemize}\setlength{\itemsep}{2pt}
  \item $T_{\min}(n)=n-1$: snabbaste indatat (redan sorterat).
  \item $T_{\max}(n)=\binom n2$: långsammaste (omvänt sorterat).
  \item $T_{\mathrm{medel}}(n)\approx n^2/4$: genomsnittet över alla ordningar.
\end{itemize}
Sedan beskriver $O$, $\Omega$ och $\Theta$ \emph{den} funktionen.
\end{column}
\end{columns}
\vfill
\begin{ahabox}
”Värsta fallet är $\Omega(n^2)$” och ”bästa fallet är $O(n)$” är helt vanliga, sanna meningar.
\end{ahabox}
\varstrip{Separation}{algoritmen}{fallet}{att fallet väljer funktionen och $O$/$\Omega$ begränsar den}
\end{frame}

\begin{frame}{En tredje axel: algoritm eller problem?}
\footnotesize
\renewcommand{\arraystretch}{1.15}
\begin{tabularx}{\textwidth}{@{}p{0.15\textwidth}XX@{}}
\toprule
& \textbf{Övre gräns} & \textbf{Undre gräns}\\
\midrule
\textbf{Algoritm $A$} & $A$ tar högst $f(n)$ på \emph{varje} indata av storlek $n$.\newline
  \textcolor{black!60}{Binärsökning: $O(\log n)$.} &
  \emph{Något} indata får $A$ att ta minst $g(n)$.\newline
  \textcolor{black!60}{Linjärsökning: $\Omega(n)$ i värsta fallet.}\\
\textbf{Problem $P$} & \emph{Det finns} en algoritm som klarar $f(n)$: $\exists A\ \forall x$.\newline
  \textcolor{black!60}{Sökning: $O(\log n)$, tack vare binärsökning.} &
  \emph{Varje} algoritm har ett indata som kostar $g(n)$: $\forall A\ \exists x$.\newline
  \textcolor{black!60}{Sökning: $\Omega(\log n)$ — uppgift 5.}\\
\bottomrule
\end{tabularx}
\medskip

\small
Möts gränserna är problemet $\Theta(\log n)$ och binärsökning optimal. En undre gräns för
problemet måste gälla även algoritmer som ingen har skrivit än.
\vfill
\begin{missbox}
”Linjärsökning tar $n$ steg, alltså är sökning $\Omega(n)$.” En gräns för en dålig
algoritm, inte för problemet.
\end{missbox}
\varstrip{Separation}{sökning i sorterad array}{algoritm eller problem}{kvantorerna $\exists A$ och $\forall A$}
\end{frame}

\begin{frame}{Uppgift 5: Undre gräns för sökning i sorterad array}
\begin{infobox}
Binärsökning i en sorterad array med $n$ tal tar som bekant tiden $O(\log n)$. Bevisa att
$\Omega(\log n)$ också är en undre gräns för antalet jämförelser som krävs för detta problem
(i värsta fallet).
\end{infobox}
\begin{columns}[T]
\begin{column}{0.44\textwidth}
\small
\textbf{Modell:} algoritmen får bara jämföra $x$ med arrayelement: $x<a_i$, $x=a_i$ eller
$x>a_i$. Varje jämförelse kostar 1, allt annat är gratis.\par\smallskip
\textbf{Beslutsträd:} en inre nod per jämförelse, ett löv per svar. En körning är en stig från
roten till ett löv. Värsta fallet är trädets höjd.
\end{column}
\begin{column}{0.53\textwidth}
\centering
{\footnotesize\bfseries\color{kthnavy}Binärsökningens träd, $n=7$}\par\smallskip
\input{gen/dt-bin7-s}\par
{\scriptsize\textcolor{kthgreen}{gröna löv}: hittad på plats $i$\quad
\textcolor{black!50}{grå löv}: finns inte}
\end{column}
\end{columns}
\end{frame}

\begin{frame}{Beviset: räkna löven}
\begin{columns}[T]
\begin{column}{0.55\textwidth}
\small
\begin{enumerate}\setlength{\itemsep}{3pt}
  \item Möjliga svar: ”$x=a_i$” för $i=1,\dots,n$, eller ”finns inte”. Minst $n+1$ olika svar,
        alltså minst $n+1$ löv.
  \item Varje inre nod har högst 3 barn. Ett träd med höjd $h$ har högst $3^h$ löv.
  \item $3^h\ge n+1\ \Rightarrow\ h\ge\log_3(n+1)=\Omega(\log n)$.
\end{enumerate}
\medskip
Det gäller \emph{varje} beslutsträd, alltså varje jämförelsebaserad algoritm. Ingen balans
behöver antas.
\end{column}
\begin{column}{0.42\textwidth}
\begin{ahabox}[fontupper=\footnotesize]
\textbf{Skarpare:} svaret ”$x=a_i$” kräver en nod som jämför med $a_i$ och får $=$.
Alltså minst $n$ inre noder. Utan $=$-löven är det ett binärträd med minst $n$ noder, och
dess höjd är minst $\lceil\log_2(n+1)\rceil$.\par\smallskip
Binärsökning gör exakt så många jämförelser. Den är optimal — inte bara asymptotiskt.
\end{ahabox}
\end{column}
\end{columns}
\vfill
\varstrip{Generalisering}{argumentet: räkna svar, jämför med höjd}{algoritmen}{att gränsen gäller alla träd på en gång}
\end{frame}

{\kaninbg
\begin{frame}{\kh Samma bevis, berättat som en motståndare}
\small
Motståndaren bestämmer inte arrayen i förväg. Den håller reda på vilka utfall som fortfarande
är möjliga: $n$ träffar och $n+1$ luckor, $2n+1$ stycken.
\begin{itemize}\setlength{\itemsep}{3pt}
  \item Vid varje jämförelse $x$ mot $a_i$ svarar den $<$ eller $>$, åt det håll där flest
        utfall finns kvar. Aldrig $=$ så länge den kan undvika det.
  \item Från $m$ möjliga utfall finns minst $(m-1)/2$ kvar efter frågan.
  \item Algoritmen kan inte svara förrän ett enda utfall återstår. Räkningen ger
        $2^k\ge n+1$, alltså minst $\lceil\log_2(n+1)\rceil$ jämförelser.
\end{itemize}
\vfill
\begin{kaninbox}[fontupper=\footnotesize]
Motståndaren \emph{är} beslutsträdet: den väljer den längsta stigen genom algoritmens träd. Samma
argument, två berättelser. Uppgift 6 använder den andra.
\end{kaninbox}
\varstrip{Fusion}{påståendet $\Omega(\log n)$}{bevistekniken: räkna löv eller spela motståndare}{att en motståndare är den värsta stigen i trädet}
\end{frame}}

\begin{frame}{Beslutsträd är inte rekursionsträd}
\begin{columns}[T]
\begin{column}{0.49\textwidth}
\centering
{\small\bfseries\color{kthnavy}Beslutsträd (binärsökning)}\par\smallskip
\input{gen/dt-bin7-s}\par\smallskip
\footnotesize En körning = \emph{en stig}. Värsta fallet = höjden.
\end{column}
\begin{column}{0.49\textwidth}
\centering
{\small\bfseries\color{kthnavy}Rekursionsträd (mergesort, $n=8$)}\par\smallskip
\begin{tikzpicture}[x=0.52cm,y=0.62cm,
  rb/.style={draw=kthblue,fill=kthlight,rounded corners=1.5pt,inner sep=1.5pt,font=\tiny,minimum height=3.2mm}]
  \node[rb,minimum width=30mm] (r) at (3.5,0) {8};
  \foreach \i in {0,1} \node[rb,minimum width=14mm] (a\i) at (1.5+4*\i,-1) {4};
  \foreach \i in {0,...,3} \node[rb,minimum width=6mm] (b\i) at (0.5+2*\i,-2) {2};
  \foreach \i in {0,...,7} \node[rb,minimum width=2.5mm] (c\i) at (\i,-3) {1};
  \foreach \i in {0,1} \draw[black!45] (r) -- (a\i);
  \foreach \i in {0,...,3} {\pgfmathtruncatemacro{\p}{\i/2}\draw[black!45] (a\p) -- (b\i);}
  \foreach \i in {0,...,7} {\pgfmathtruncatemacro{\p}{\i/2}\draw[black!45] (b\p) -- (c\i);}
\end{tikzpicture}\par\smallskip
\footnotesize En körning = \emph{hela trädet}. Tid = summan.
\end{column}
\end{columns}
\vfill
\begin{missbox}
”Binärsökningens rekursion har djup $\log n$, alltså är $\Omega(\log n)$ bevisat.” Det säger vad
\emph{binärsökning} gör, inte vad alla algoritmer måste göra.
\end{missbox}
\varstrip{Kontrast}{”träd som beskriver en algoritm”}{vad en körning är: stig eller hela trädet}{att beslutsträdet beskriver alla körningar}
\end{frame}

\begin{frame}{Två algoritmer, två träd, lika många löv}
\begin{columns}[T]
\begin{column}{0.4\textwidth}
\centering
{\small\bfseries\color{kthnavy}Linjärsökning: höjd 7}\par\smallskip
\begin{tikzpicture}[scale=0.6,dn/.style={circle,draw=kthnavy,fill=white,line width=0.6pt,inner sep=0pt,minimum size=6.2mm,font=\scriptsize},
  hit/.style={rectangle,draw=kthgreen,fill=kthmint,inner sep=0pt,minimum size=3.6mm,font=\tiny\bfseries},
  ej/.style={rectangle,draw=black!35,fill=black!8,inner sep=0pt,minimum size=2.6mm},
  te/.style={draw=black!45,line width=0.6pt},
  tl/.style={font=\tiny,text=black!60,inner sep=0.5pt,fill=white}]
  \node[dn] (n1) at (0.0,-0.0) {$a_1$};
  \node[ej] (el1) at (-0.62,-0.72) {};
  \node[hit] (h1) at (0.0,-0.72) {1};
  \draw[te] (n1) -- (el1); \draw[te] (n1) -- (h1);
  \node[dn] (n2) at (1.25,-0.72) {$a_2$};
  \node[ej] (el2) at (0.63,-1.44) {};
  \node[hit] (h2) at (1.25,-1.44) {2};
  \draw[te] (n2) -- (el2); \draw[te] (n2) -- (h2);
  \draw[te] (n1) -- (n2);
  \node[dn] (n3) at (2.5,-1.44) {$a_3$};
  \node[ej] (el3) at (1.88,-2.16) {};
  \node[hit] (h3) at (2.5,-2.16) {3};
  \draw[te] (n3) -- (el3); \draw[te] (n3) -- (h3);
  \draw[te] (n2) -- (n3);
  \node[dn] (n4) at (3.75,-2.16) {$a_4$};
  \node[ej] (el4) at (3.13,-2.88) {};
  \node[hit] (h4) at (3.75,-2.88) {4};
  \draw[te] (n4) -- (el4); \draw[te] (n4) -- (h4);
  \draw[te] (n3) -- (n4);
  \node[dn] (n5) at (5.0,-2.88) {$a_5$};
  \node[ej] (el5) at (4.38,-3.5999999999999996) {};
  \node[hit] (h5) at (5.0,-3.5999999999999996) {5};
  \draw[te] (n5) -- (el5); \draw[te] (n5) -- (h5);
  \draw[te] (n4) -- (n5);
  \node[dn] (n6) at (6.25,-3.5999999999999996) {$a_6$};
  \node[ej] (el6) at (5.63,-4.319999999999999) {};
  \node[hit] (h6) at (6.25,-4.319999999999999) {6};
  \draw[te] (n6) -- (el6); \draw[te] (n6) -- (h6);
  \draw[te] (n5) -- (n6);
  \node[dn] (n7) at (7.5,-4.32) {$a_7$};
  \node[ej] (el7) at (6.88,-5.04) {};
  \node[hit] (h7) at (7.5,-5.04) {7};
  \draw[te] (n7) -- (el7); \draw[te] (n7) -- (h7);
  \draw[te] (n6) -- (n7);
  \node[ej] (last) at (8.12,-5.04) {};
  \draw[te] (n7) -- (last);
\end{tikzpicture}

\end{column}
\begin{column}{0.57\textwidth}
\centering
{\small\bfseries\color{kthnavy}Binärsökning: höjd 3}\par\smallskip
\begin{tikzpicture}[scale=0.7,dn/.style={circle,draw=kthnavy,fill=white,line width=0.6pt,inner sep=0pt,minimum size=6.2mm,font=\scriptsize},
  hit/.style={rectangle,draw=kthgreen,fill=kthmint,inner sep=0pt,minimum size=3.6mm,font=\tiny\bfseries},
  ej/.style={rectangle,draw=black!35,fill=black!8,inner sep=0pt,minimum size=2.6mm},
  te/.style={draw=black!45,line width=0.6pt},
  tl/.style={font=\tiny,text=black!60,inner sep=0.5pt,fill=white}]
  \node[dn] (n4) at (0,0) {$a_4$};
  \node[dn] (n2) at (-3.4,-1.1) {$a_2$};
  \node[dn] (n6) at (3.4,-1.1) {$a_6$};
  \node[dn] (n1) at (-5.1,-2.2) {$a_1$};
  \node[dn] (n3) at (-1.7,-2.2) {$a_3$};
  \node[dn] (n5) at (1.7,-2.2) {$a_5$};
  \node[dn] (n7) at (5.1,-2.2) {$a_7$};
  \draw[te] (n4) -- node[tl,pos=0.45] {$<$} (n2);
  \draw[te] (n4) -- node[tl,pos=0.45] {$>$} (n6);
  \node[hit] (h4) at (0,-1.1) {4};
  \draw[te] (n4) -- node[tl,pos=0.5] {$=$} (h4);
  \draw[te] (n2) -- node[tl,pos=0.45] {$<$} (n1);
  \draw[te] (n2) -- node[tl,pos=0.45] {$>$} (n3);
  \node[hit] (h2) at (-3.4,-2.2) {2};
  \draw[te] (n2) -- node[tl,pos=0.5] {$=$} (h2);
  \draw[te] (n6) -- node[tl,pos=0.45] {$<$} (n5);
  \draw[te] (n6) -- node[tl,pos=0.45] {$>$} (n7);
  \node[hit] (h6) at (3.4,-2.2) {6};
  \draw[te] (n6) -- node[tl,pos=0.5] {$=$} (h6);
  \node[ej] (el1) at (-5.75,-3.25) {};
  \node[hit] (h1) at (-5.1,-3.3000000000000003) {1};
  \node[ej] (er1) at (-4.449999999999999,-3.25) {};
  \draw[te] (n1) -- (el1); \draw[te] (n1) -- (h1); \draw[te] (n1) -- (er1);
  \node[ej] (el3) at (-2.35,-3.25) {};
  \node[hit] (h3) at (-1.7,-3.3000000000000003) {3};
  \node[ej] (er3) at (-1.0499999999999998,-3.25) {};
  \draw[te] (n3) -- (el3); \draw[te] (n3) -- (h3); \draw[te] (n3) -- (er3);
  \node[ej] (el5) at (1.0499999999999998,-3.25) {};
  \node[hit] (h5) at (1.7,-3.3000000000000003) {5};
  \node[ej] (er5) at (2.35,-3.25) {};
  \draw[te] (n5) -- (el5); \draw[te] (n5) -- (h5); \draw[te] (n5) -- (er5);
  \node[ej] (el7) at (4.449999999999999,-3.25) {};
  \node[hit] (h7) at (5.1,-3.3000000000000003) {7};
  \node[ej] (er7) at (5.75,-3.25) {};
  \draw[te] (n7) -- (el7); \draw[te] (n7) -- (h7); \draw[te] (n7) -- (er7);
\end{tikzpicture}

\end{column}
\end{columns}
\vfill
\small
Båda har $7$ gröna och $8$ grå löv. Gränsen säger inte att alla träd är balanserade. Den säger
att \emph{inget} träd med så många löv kan vara lägre än $\lceil\log_2 8\rceil=3$.
\varstrip{Generalisering}{problemet och antalet löv}{algoritmen, alltså trädets form}{att den undre gränsen gäller varje form}
\end{frame}

\begin{frame}{Fyra verktyg för undre gränser}
\begin{columns}[T]
\begin{column}{0.68\textwidth}
\footnotesize
\renewcommand{\arraystretch}{1.35}
\begin{tabularx}{\linewidth}{@{}p{0.2\linewidth}Xp{0.27\linewidth}@{}}
\toprule
\textbf{Verktyg} & \textbf{Idé} & \textbf{Idag}\\
\midrule
Läsa och skriva & Svaret har storlek $s$ $\Rightarrow$ $\Omega(s)$ & Euler: $\Omega(|E|)$\\
Beslutsträd & $L$ möjliga svar $\Rightarrow$ höjd $\ge\log L$ & Uppgift 5; sortering\\
Motståndare & Svara elakt men konsekvent & Uppgift 5 igen; uppgift 6\\
Reduktion & Löser du $B$ snabbt, löser du $A$ snabbt & Sortering $\to$ konvext hölje\\
\bottomrule
\end{tabularx}
\end{column}
\begin{column}{0.29\textwidth}
\centering
\begin{tikzpicture}[x=0.75cm,y=0.42cm]
  \draw[->,black!50] (-2.3,0) -- (2.4,0) node[right,font=\tiny] {$x$};
  \draw[->,black!50] (0,-0.2) -- (0,4.2) node[above,font=\tiny] {$x^2$};
  \draw[black!25,domain=-2.0:2.05,samples=40] plot (\x,{\x*\x});
  \draw[kthblue,line width=1pt] (-1.8,3.240) -- (-0.7,0.490) -- (0.4,0.160) -- (1.1,1.210) -- (1.9,3.610) -- cycle;
  \fill[kthorange] (-1.8,3.240) circle (2pt);
  \fill[kthorange] (-0.7,0.490) circle (2pt);
  \fill[kthorange] (0.4,0.160) circle (2pt);
  \fill[kthorange] (1.1,1.210) circle (2pt);
  \fill[kthorange] (1.9,3.610) circle (2pt);
\end{tikzpicture}
\par
\scriptsize Sortera $x_1,\dots,x_n$: beräkna konvexa höljet av punkterna $(x_i,x_i^2)$. Höljet
går igenom dem i sorterad ordning.
\end{column}
\end{columns}
\vfill
\begin{ahabox}
Reduktionen går \emph{från} det kända problemet \emph{till} det nya: om konvexa höljet gick
snabbare än $\Omega(n\log n)$ skulle sortering också göra det.
\end{ahabox}
\end{frame}

{\kaninbg
\begin{frame}{\kh Sortering: $n!$ löv}
\small
\begin{itemize}\setlength{\itemsep}{3pt}
  \item En jämförelsebaserad sortering måste kunna ge var och en av de $n!$ ordningarna som svar.
  \item Jämförelsen $a_i<a_j$ har två utfall. Ett binärt träd med höjd $h$ har högst $2^h$ löv.
  \item $2^h\ge n!\ \Rightarrow\ h\ge\log_2 n!\ \ge\ \frac n2\log_2\frac n2=\Omega(n\log n)$,
        eftersom $n!$ har minst $n/2$ faktorer som är $\ge n/2$.
  \item Mergesort klarar $O(n\log n)$. Problemet är alltså $\Theta(n\log n)$.
\end{itemize}
\vfill
\begin{kaninbox}[fontupper=\footnotesize]
Samma recept som i uppgift 5, fler löv. För $n=7$: $\lceil\log_2 5040\rceil=13$ jämförelser
krävs i värsta fallet. Det finns algoritmer som klarar sig med exakt 13, men ingen känd
formel för det bästa antalet i allmänhet.
\end{kaninbox}
\end{frame}}

{\kaninbg
\begin{frame}{\kh Modellen spelar roll}
\small
Undre gränsen gäller algoritmer som bara \emph{jämför}. Den som räknar med nycklarna har inte
lovat något:
\begin{itemize}\setlength{\itemsep}{3pt}
  \item \textbf{Interpolationssökning} gissar var $x$ borde ligga utifrån värdet. Jämnt
        fördelade nycklar: $O(\log\log n)$ i medel. Värsta fallet är fortfarande $\Theta(n)$.
  \item \textbf{Hashning}: $O(1)$ förväntat. Behöver inte ens sorterad ordning.
  \item \textbf{Räknesortering, radixsortering}: $O(n+k)$. ”Sortering är $\Omega(n\log n)$”
        gäller jämförelser.
\end{itemize}
\vfill
\begin{missbox}[fontupper=\footnotesize]
”Hashning motbevisar uppgift 5.” Nej: uppgift 5 handlar om jämförelser. En undre gräns är
alltid en gräns \emph{i en modell}.
\end{missbox}
\varstrip{Kontrast}{problemet: hitta $x$}{beräkningsmodellen}{att en undre gräns bara gäller i sin modell}
\end{frame}}

\begin{frame}{Vanliga fel i uppgift 5}
\begin{columns}[T]
\begin{column}{0.49\textwidth}
\begin{missbox}[fontupper=\footnotesize]
”Trädet har $n$ löv.” Svaret ”finns inte” är också ett svar: minst $n+1$.
\end{missbox}
\begin{missbox}[fontupper=\footnotesize]
”Höjden är $O(\log n)$.” Det gäller binärsökningens träd. Beviset behöver att \emph{varje}
träd har höjd $\Omega(\log n)$.
\end{missbox}
\end{column}
\begin{column}{0.49\textwidth}
\begin{missbox}[fontupper=\footnotesize]
”Trädet är balanserat, så höjden är $\log n$.” Balans får man inte anta: trädet tillhör en
okänd algoritm. Ett träd med $L$ löv har höjd $\ge\log L$ ändå.
\end{missbox}
\begin{missbox}[fontupper=\footnotesize]
”Jag hittade ett indata där min algoritm är långsam.” Det är en undre gräns för \emph{din}
algoritm. Problemet kräver $\forall A$.
\end{missbox}
\end{column}
\end{columns}
\varstrip{Kontrast}{uppgiften}{beviset, rätt och nästan rätt}{löv, riktning, balans, kvantor}
\end{frame}
